The classical statement allows an arbitrary function $\psi:\mathbb N\to[0,\infty)$:
\[
 \left|\alpha-\frac aq\right|<\frac{\psi(q)}q,\quad (a,q)=1
 \tag{1}
\]
has infinitely many solutions for almost every $\alpha$ if and only if
$\sum_q\phi(q)\psi(q)/q=\infty$.
The current page prints $\mathbb N\to\mathbb N$. If $\mathbb N$ means positive integers in that display, then $\psi\ge1$, the sum dominates the harmonic series, and continued-fraction convergents already give infinitely many solutions for every irrational $\alpha$. We address the intended real-valued statement instead.

Koukoulopoulos and Maynard proved it in \emph{On the Duffin--Schaeffer conjecture}, Ann. of Math. 192 (2020), 251--307.
The proof below reconstructs the reduction from their weighted exceptional-pair estimate, including the summation over exceptional scales. That deep estimate and the classical auxiliary results are explicitly imported; their full GCD-graph proof is not reproduced. No novelty is claimed.
